\section{Introduction}
\label{sec:introduction}

When something must be found quickly over a large area --- a missing hiker,
survivors in a collapsed structure, the front of a spreading wildfire ---
teams of small autonomous aircraft are increasingly the tool of choice: they
are cheap, expendable, fast to deploy, and, above all, \emph{parallel}
\citep{queralta2020collaborative,julian2019distributed}. Yet the hard
problem in multi-drone search is not flight; it is coordination under
scarcity. Each vehicle perceives only a small neighbourhood of a large,
unknown environment; communication reaches only nearby teammates, so no
shared global picture exists; and the mission clock is unforgiving --- in
search and rescue, time \emph{is} the objective. A team that cannot divide
the search space effectively degenerates into an expensive single searcher.

This project adds to that already delicate setting the ingredient that
real operations always contain and idealized models usually omit:
\textbf{failure}. Small drones break --- motors fail, batteries die,
weather intervenes --- and they break \emph{mid-mission}, unpredictably,
without ceremony. A search system that presumes a fixed team is brittle in
exactly the situations that matter most. Worse, in a genuinely
decentralized setting a drone's death is not even \emph{announced}: its
teammates must discover the loss through their own sensing and
communication, and until they do, they coordinate with a phantom. Fault
tolerance in a multi-agent system is thus not one problem but three ---
surviving the loss (redundancy), learning of it (decentralized fault
knowledge), and adapting to it (re-planning the division of labour) --- and
this report treats all three as first-class design objects.

\paragraph{The problem.}
Concretely, we study the following task, formalized in
Section~\ref{sec:problem}. A team of $\nagents$ drones spawns co-located on
a corner of a $\gridsize\times\gridsize$ occupancy grid scattered with
obstacles; a single target, guaranteed reachable, is hidden at an unknown
cell. Each drone sees only a window of radius $\rvis$ around itself,
remembers what it has seen, and fuses its accumulated map with any teammate
within communication range $\rcomm$. The team has a hard budget of
$T_{\max}$ turns to physically reach the target. At every turn, any live
drone may irreversibly break with probability $\pfault$; a wreck stops
sensing, moving, and communicating, its turns still consume the budget, and
no one is notified --- fault knowledge spreads only through a passive
distress beacon and the ordinary map fusion. Success is binary and
team-level: \emph{some} drone reaches the target in time, or the episode
fails.

\paragraph{The approach.}
We solve the task with cooperative deep reinforcement learning under
parameter sharing \citep{gupta2017cooperative}: a single Dueling Double DQN
\citep{mnih2015human,vanhasselt2016double,wang2016dueling} drives every
drone, consuming that drone's private maps through a convolutional
architecture with an attention bottleneck, and trained on the pooled
experience of the whole team. Rather than fixing one operating condition,
every training episode randomizes the regime --- vision radius,
communication range, team size, obstacle density, and fault probability
\citep{tobin2017domain} --- and the sampled parameters that a real drone
could plausibly know are fed back to the network as a context vector,
together with two per-drone quantities: an identity (which breaks the
symmetry of parameter sharing and lets co-located drones adopt different
roles from the first step) and the drone's own \emph{belief} about how many
teammates remain operational. The result is a single generalist policy
that searches alone or in a team, myopic or far-sighted, in empty fields or
cluttered mazes --- and that experiences teammate death, and the need to
absorb it, as an ordinary feature of its world. The fault model demands
equal care on the learning side: a fault is exogenous bad luck, not a
consequence of behaviour, and Section~\ref{sec:training} shows why treating
it as anything but a truncation would silently bias the learned values.

\paragraph{Why this is a multi-agent systems study.}
A reinforcement-learning pipeline, however carefully engineered, is not by
itself the point of this report. The system is deliberately read, in
Section~\ref{sec:mas}, through the conceptual apparatus of multi-agent
systems \citep{wooldridge2009introduction,jennings2001agent}: the drones
are audited against the defining properties of agenthood; the environment
is analysed as a first-class abstraction that mediates and governs all
interaction; the map fusion is read as knowledge-level, protocol-free
communication over a distributed shared dataspace; the visited-cell traces
as a virtualized form of stigmergy; the fused map, the beacon, and the
action mask as coordination artefacts; the team's dispersion and post-fault
re-covering as emergent, self-organized structure that no component
specifies; and the whole design is placed on the autonomy spectrum ---
operationally autonomous, motivationally not. This analysis is not
decorative: it is what connects a concrete learned system to the questions
--- about coordination, openness, resilience, and responsibility --- that
the multi-agent perspective exists to ask. The same lens then carries the
final part of the report, where autonomy without deliberation and
competence without transparency become ethical and legal problems rather
than engineering trade-offs.

\paragraph{Contributions.}
The report makes the following contributions:

\begin{enumerate}[leftmargin=2em]
  \item \emph{A cooperative search environment with a decentralized fault
    model} (Sections~\ref{sec:problem}--\ref{sec:architecture}): a
    configurable grid world, built from scratch, in which drone failures
    are permanent, unannounced, and discoverable only through a distress
    beacon and range-limited map fusion --- no component ever accesses a
    global oracle of team state.
  \item \emph{A single generalist policy via domain randomization and
    context conditioning} (Sections~\ref{sec:architecture}
    and~\ref{sec:training}): one shared network that generalizes across
    team sizes, sensing and communication regimes, obstacle densities, and
    fault rates, with per-drone identity and believed-alive-count inputs
    enabling role differentiation and post-fault adaptation.
  \item \emph{A set of reinforcement-learning correctness decisions for
    fault-prone teams} (Section~\ref{subsec:train-correctness}): faults as
    bootstrapped truncations, retroactive team credit restricted to
    surviving drones, and shaping-robust checkpoint selection --- each with
    an explicit argument for why the obvious alternative is subtly wrong.
  \item \emph{A systematic multi-agent reading of the system}
    (Section~\ref{sec:mas}): a detailed mapping between the implemented
    mechanisms and the paradigm's core concepts --- agenthood, environment,
    coordination, artefacts, self-organisation, openness, autonomy, and
    multi-agent learning.
  \item \emph{An experimental protocol --- and, pending the conclusion of
    training, its results} (Section~\ref{sec:experiments}): seeded,
    reproducible evaluation of fault robustness, generalization across
    randomization axes, and the quantitative signatures of emergent
    coordination.
  \item \emph{A prospective ethical and regulatory analysis}
    (Section~\ref{sec:ethics}): the capability's dual-use profile, its
    treatment under the GDPR and the EU AI Act (including the military
    carve-out that exempts precisely the gravest use), an argument that
    the system as designed could not lawfully or ethically serve as an
    autonomous weapon, and a set of design commitments that follow.
\end{enumerate}

\paragraph{Roadmap.}
Section~\ref{sec:background} reviews the reinforcement-learning and
multi-agent-learning toolbox. Section~\ref{sec:problem} formalizes the task
as a partially observable, identical-interest Markov game with sequential
execution and defines the observation, communication, and fault models.
Section~\ref{sec:architecture} describes the implementation: environment,
observation encoding, neural architecture, and learning machinery.
Section~\ref{sec:training} covers the training methodology and its
correctness decisions. Section~\ref{sec:mas} develops the multi-agent
analysis at length. Section~\ref{sec:experiments} presents the experimental
design and results. Section~\ref{sec:ethics} develops the ethical and
regulatory analysis, and Section~\ref{sec:conclusions} concludes.
Appendices consolidate configuration details and reproduction instructions.

% TODO (figures pass): GUI teaser screenshot (4 drones, one wreck emitting
% a beacon) once final visuals are ready.
